\section{Discussion}\label{sec:discussion}

\paragraph{What the paper claims, and what it does not.} The conceptual contribution is the contradiction triangle (Theorem~\ref{thm:contradiction}): no single input can be moved far enough to rescue an infeasible claim without the move itself being visible against that input's independent measurement, and the contradiction radius $\rho$ gives a unit-free measure of how far a claim sits from the feasible set $\mathcal F$. Whereas \citet{spagat2009} frame war-death estimation as an \emph{arena of contestation}, our framework gives the analyst a numerical referee for that arena, in the spirit of \citet{lumpricebanks2013} for multiple systems estimation and \citet{radford2023} for Bayesian source aggregation. In the Gaza application the referee's ruling is deliberately modest: the data bound the combatant share of direct deaths at $\nQOneRound{}\%$ under only the minimal exposure assumption $\mu\ge 1$ and at about $\nQTwoRound{}\%$ under historically calibrated exposure, and the public claim of $\nIdfQLoRound{}$--$\nIdfQHiRound{}\%$ is rejected under every anchor and every calibrated exposure assumption we examined---its upper endpoint is arithmetically infeasible for any $\mu\ge 1$, and its lower endpoint survives only by asserting that civilian men were no more exposed than women and children ($\mu\approx 1$), against the calibration evidence---and much of the tension is internal, since the claimant's own pre-war stock estimate, the toll it accepts, and its intelligence database already contradict its public claim before any external measurement is consulted. The paper does \emph{not} claim that the spatial model's $q=\nSpatialQ{}\%$ is the truth; that number is one prior-dependent point inside the set, reported with its full model documentation so that readers can locate exactly which assumptions produce it.

\paragraph{Where the power comes from.} It is worth being explicit about why the bound binds. The demographic relation has one free behavioural parameter, the exposure ratio $\mu$, and the data pin the product, not the factors: a given deficit of women and children among the dead can be explained by combatants dying, by civilian men being more exposed, or by a mixture. High-combatant claims must therefore assert that $\mu\approx 1$---that civilian men in an urban war were no more exposed than women and children---at the same time as independent calibrations on conflicts with recorded combatant status put $\mu$ at $2$ or higher \citep{frost2026,cockerill2024}. The claim is not rejected by any single measurement but by the joint configuration; that is the under-fudgibility mechanism doing its work, and it is also why the diagnostic is hard to game: an actor wishing to defeat it must either assert a large, visible error at a single institution or coordinate distortions across a census bureau, a health ministry's death records, foreign military-balance assessments, and historical calibration data that predate the conflict.

The same structure prices the common defences of the claim. Each rescues feasibility on one axis by conceding ground elsewhere: large wartime recruitment relaxes the manpower constraint, but concedes that the war generates fighters faster than it removes them; differential undercounting of combatant deaths in the MoH record relaxes the demographic constraint, but concedes that the headline death toll is too low; a much larger true toll relaxes the arithmetic, but concedes a war more lethal than the claim's proponents otherwise allow. The framework does not adjudicate among these defences; it computes what each costs in standardised units of the input it moves (Section~\ref{sec:contradiction}) and forces the trade-off into the open.

\paragraph{Relation to convergent evidence.} The convergence table (Table~\ref{tab:convergence}) is, to our knowledge, the first side-by-side placement of the demographic-decomposition literature, the named-combatant documentary record, and a generative spatial model on the same scale. We read the residual disagreement---roughly $q\in[\nSpatialQRound{}\%,\nQOneCeil{}\%]$---as an honest statement of what current public data cannot resolve: it is essentially the width of the exposure-calibration question plus the classification question (who counts as a combatant, on which the named-list evidence and demographic methods can legitimately differ). Narrowing it further requires either a defensible conflict-specific measurement of $\mu$ or transparent, auditable combatant lists, not another aggregation rule.

\paragraph{Limitations.}
\begin{enumerate}[leftmargin=*]
\item Class $\mathrm{AM}$ is treated as a single bucket; finer age partitions (18--29, 30--49, 50$+$) tighten the bounds when age-resolved data exist, and the calibration of $\mu$ from near-cutoff age bands \citep{frost2026} is itself subject to sampling noise in small historical samples.
\item The exposure range $[\mu_{\mathrm{lo}},\mu_{\mathrm{hi}}]$ is imported from historical calibration; if the present conflict's exposure structure is genuinely outside that range, only the exposure-agnostic bound and the manpower bound survive. We report both throughout for this reason.
\item The radius standardises by a stated uncertainty scale (in the application, a conservative multiple of the sampling error); systematic error in an anchor (e.g.\ selection in who gets recorded as dead) is addressed by re-computing under alternative anchors, not by the radius itself. Section~\ref{sec:gaza} therefore reports all anchors, with the most claim-favourable one as primary.
\item The radius makes no probabilistic claim about how likely a given number of standardised units is; where a defensible sampling distribution exists, a bootstrap or non-parametric posterior is preferred for small demographic samples.
\item Combatant status is a legal-institutional category, not a demographic one. The framework bounds the share of deaths attributable to members of armed groups as counted by manpower estimates; disputes about who should count as a combatant move $M$ and the interpretation of named lists, and enter the diagnostic through those channels.
\item The estimand is the combatant share of \emph{direct} deaths throughout; indirect deaths from blockade, famine, or healthcare collapse require separate excess-mortality methodology \citep{checchiroberts2008,burnham2006,degommeguhasapir2010,gomez2025} and are outside the frame. The denominator scenarios varied in Section~\ref{sec:gaza} (survey-corrected and missing-persons totals) are corrections to the count of \emph{direct} deaths, so they change $D$ without changing the estimand.
\item The combatant stock $M$ and population shares $(w,a,f)$ are treated as fixed over the analysis window: wartime recruitment, ageing into the adult-male class, and displacement out of the territory are not modelled. Recruitment is the quantitatively relevant channel and is priced explicitly by the manpower ledger of Section~\ref{sec:gaza}; the demographic drifts are plausibly second-order at this window length.
\end{enumerate}

\paragraph{Reproducibility.} All bounds, radii, and derived table entries in this paper are computed by validation scripts that recompute them from the raw inputs, emit the numeric macros and tables the text includes, and fail on any internal discrepancy. In addition, the algebraic core of the framework (Appendix~\ref{app:proofs}) is formalised and machine-checked in Lean~4 with mathlib, and the remaining results are checked symbolically and exercised numerically by dedicated scripts. Code, the formal proofs, the spatial simulator, and all inputs are available in the companion repository.
